\section{QA · LaTeX note}
% #qa/latex

A note written in LaTeX. Check: each construct is drawn as its heading says, its LaTeX shows while the cursor is on it, and the title in the list is "QA · LaTeX note".

\paragraph{1. Headings}

\subsection{A subsection}
\subsubsection*{A starred subsubsection}

\paragraph{2. Styles}

\textbf{bold}, \textit{italic}, \emph{emphasis}, \texttt{code}, \underline{underlined}, \textsc{Small Caps}, \textbf{bold with \emph{emphasis} inside}.

\paragraph{3. Characters}

50\% of \$10 \& \#1, a~tie, en--dash, em---dash, ``quoted'', and so on\ldots{} \LaTeX.

\paragraph{4. Lists}

\begin{itemize}
  \item First
  \item Second
  \begin{enumerate}
    \item Nested, numbered
    \item Again
  \end{enumerate}
  \item[--] Its own marker
\end{itemize}

\begin{description}
  \item[Term] Its description.
\end{description}

\paragraph{5. Math}

Inline $e^{i\pi} + 1 = 0$ and \(a^2 + b^2 = c^2\).

\[ \sum_{n=1}^\infty \frac{1}{n^2} = \frac{\pi^2}{6} \]

\begin{equation}
  \int_0^1 x^2\,dx = \frac{1}{3} \label{eq:integral}
\end{equation}

\begin{align*}
  a &= b + c \\
  d &= e
\end{align*}

\paragraph{6. Broken math}

Red, as written, with the reason on hover: $\frac{1}{$

\paragraph{7. Table}

\begin{tabular}{|l|c|r|}
  \hline
  Left & Center & Right \\
  \hline
  a & $x^2$ & 1 \\
  longer cell & b & 22 \\
  \hline
\end{tabular}

\paragraph{8. Quote and centered}

\begin{quote}
A quote.
\end{quote}

\begin{center}
Centered text.
\end{center}

\paragraph{9. Code}

\begin{verbatim}
raw \textbf{not bold} $not math$
\end{verbatim}

\paragraph{10. Links and footnotes}

A \href{https://example.com}{named link}, a bare \url{https://example.com}, and a footnote.\footnote{The footnote's text.}

\paragraph{11. Links drawn short}

\url{https://arxiv.org/abs/1706.03762} reads arXiv:1706.03762 and \url{https://github.com/overtio/lines/issues/12} overtio/lines\#12, each with its mark; \href{https://arxiv.org/abs/2006.11239}{a named one} keeps its text with the mark in front.

\paragraph{12. Comments and other commands}

Text % a comment with #qa/latex-tag
\maketitle \vspace{1em} \label{x} stay faint.

\paragraph{13. Line break}

First line\\
second line.

\paragraph{14. Theorems}

\newtheorem{lem}[theorem]{Lemma}

\begin{lem}[Small]
A lemma, numbered along with theorems: Lemma 1.
\end{lem}

\begin{theorem}
Theorem 2, with math in its frame:
\[ a^2 + b^2 = c^2 \]
\end{theorem}

\begin{proof}
Italic Proof., and a box at the end.
\end{proof}

\paragraph{15. Citations and references}

As shown in \cite{knuth1984, lamport1994} and Section~\ref{sec:intro}, see \eqref{eq:one}. \label{qa:hidden} The label before this is hidden.

\paragraph{16. Own math commands}

\newcommand{\R}{\mathbb{R}}
\newcommand{\abs}[1]{\left\lvert #1 \right\rvert}
\DeclareMathOperator{\sinc}{sinc}
Drawn with the note's commands: $\abs{x} \in \R$ and $\sinc x$.

\paragraph{17. Accents, letters, verb}

Gödel, Erdős, Łukasiewicz, façade: G\"odel, Erd\H{o}s, \L ukasiewicz, fa\c{c}ade. Code: \verb|x = {a & b}|.

\paragraph{18. Nested numbering}

\begin{enumerate}
  \item One
  \begin{enumerate}
    \item (a)
    \begin{enumerate}
      \item i.
    \end{enumerate}
  \end{enumerate}
\end{enumerate}

\paragraph{19. Bibliography}

\begin{thebibliography}{9}
\bibitem{knuth1984} D. Knuth, \emph{The \TeX book}, 1984.
\bibitem{lamport1994} L. Lamport, \emph{\LaTeX: A Document Preparation System}, 1994.
\end{thebibliography}

\paragraph{20. Not closed}

\textbf{never closed, and \begin{itemize} with no end: shown as written.
